\section{Introduction}\label{sec:introduction}

The Fontaine--Mazur conjecture predicts that irreducible \(p\)-adic
Galois representations satisfying suitable finiteness and local
conditions arise from algebraic geometry \cite{FM1995}. Modular forms supply
a central class of examples through Deligne's cohomological construction
\cite[Sections~3.9--3.12 and~5.2--5.5]{Deligne1971}. Over
\(\mathbb Q\), its odd two-dimensional form has a more precise
modularity formulation. A continuous representation
\[
 r:G_{\mathbb Q}\longrightarrow\GL_2(\overline{\mathbb Q}_p)
\]
is \emph{odd} if \(\det r(c)=-1\) for a complex conjugation \(c\).
It is \emph{regular de Rham} at \(p\) if its restriction to
\(G_{\mathbb Q_p}\) is de Rham and has two distinct Hodge--Tate
weights. The conjecture asserts that an irreducible such
representation, unramified outside finitely many primes, comes
from a classical cuspidal eigenform, allowing a Tate twist.

\begin{theorem}\label{thm:main}
Let
\[
 r:G_{\mathbb Q}\longrightarrow\GL_2(\overline{\mathbb Q}_2)
\]
be continuous and irreducible. Suppose that \(r\) is odd, is
unramified outside finitely many primes, and is de Rham at \(2\)
with distinct Hodge--Tate weights. Then \(r\) is isomorphic to
a Tate twist of the \(2\)-adic Galois representation associated
with a classical cuspidal eigenform.
\end{theorem}

Thus Theorem~\ref{thm:main} resolves the odd, regular two-dimensional
Fontaine--Mazur conjecture over \(\mathbb Q\) at \(2\). Its scope includes
every semisimple residual representation, including scalar and reducible
ones. Together with the odd-prime results discussed below, it gives the
stated modularity conclusion at every prime. The geometric realization
of modular representations also gives the geometric-origin conclusion
in this odd, regular setting.

Theorem~\ref{thm:main}, in particular its absence of a residual-image
hypothesis, supplies the weight-two modularity input in the height argument
of \emph{Hilbert's tenth problem over the rational numbers}
\cite[Section~4]{OpenAIH10Q2026}, whose undecidability proof also uses
substantial additional arithmetic and logical constructions.

\subsection{Background and the dyadic problem}

The deformation-theoretic approach to modularity lifting developed
by Wiles and Taylor--Wiles compares Galois deformation rings with
Hecke algebras \cite{Wiles1995,TaylorWiles1995}. Serre's conjecture
concerns the residual starting point \cite{Serre1987}; its proof by
Khare--Wintenberger and Kisin supplies modularity for odd absolutely
irreducible residual representations, including characteristic two
\cite{KhareWintenberger2009,KhareWintenberger2009II,Kisin2009}.
The existence of a residual modular form still leaves the problem
of proving modularity for a specified characteristic-zero lift.
For a reducible residual representation, even the modular seed
must be obtained by a different construction.

Skinner and Wiles developed ordinary families and pro-modularity
arguments for residually reducible representations and then for
nearly ordinary deformations of irreducible residual representations
\cite{SkinnerWiles1999,SkinnerWiles2001}. Allen adapted this family
method to nearly ordinary dyadic representations with solvable
absolutely irreducible residual image \cite{Allen2014}.
Kisin brought local deformation geometry and the $p$-adic
Langlands correspondence into the regular Fontaine--Mazur problem
\cite{Kisin2009FM}. Emerton's local--global compatibility work
separated occurrence in completed cohomology from the production
of classical vectors \cite{Emerton2011}. These developments
explain why both stages of the proof below are needed, and why
ordinary auxiliary families can be useful without ordinarity
being imposed on the final target.

At odd primes, Pan's work treats the residually reducible case and,
together with preceding modularity lifting theorems, establishes the
regular conjecture apart from one local residual configuration at~$3$
\cite[Theorems~1.0.2 and~1.0.4]{Pan2022}. Zhang removes that
configuration, obtaining the stated odd, regular theorem for every
odd prime \cite[Theorem~1.0.2]{Zhang2025}. At~$2$, Tung removes
the remaining local restriction in a modularity lifting theorem whose
global residual image is nonsolvable
\cite[Theorem~A]{Tung2021Dyadic}. Thus local scalar residual
behavior and a globally scalar residual representation must be
distinguished: a local theorem covering the former does not by
itself settle the latter.

The proof has a useful precedent in the passage from unrestricted
pro-modularity to classicality developed by Pan, and in Zhang's
potential pro-modularity argument \cite{Pan2022,Zhang2025}.
The present dyadic problem requires integral control of the local
actions and of deformation families even when their residual
representation is scalar. We first construct potential occurrence
in completed definite forms, then extract classical vectors, and
finally descend. The distinction between occurrence and classicality
also appears in Pan's theorem for representations in the completed
cohomology of modular curves: that theorem assumes occurrence as
a hypothesis \cite[Theorem~7.1.2]{Pan2026Classicality}.

Thorne's theorem supplies an essential set of classical points:
ordinary, potentially crystalline representations of weights
$\{0,1\}$ are modular at every prime, without a residual
irreducibility hypothesis \cite[Theorem~D]{Thorne2026}.
After a global normalization and passage to the dual, our family
construction makes this theorem applicable at a dense
set of points, and the subsequent specialization argument recovers
the original target's arbitrary weight gap. Prescribing global
lifts with chosen local components, as in Emerton--Gee--Pan--Zhu
\cite[Theorem~1.1.1]{EGPZ2026}, is a different conclusion;
their dyadic theorem also retains a nonsolvable residual-image
hypothesis.

The local input is the \(p\)-adic Langlands correspondence of
Colmez~\cite{Colmez2010}, in the all-prime form of Colmez, Dospinescu
and Pa\v{s}k\=unas~\cite{CDP2014}. Pa\v{s}k\=unas and Tung's block
finiteness theorem~\cite{PaskunasTung2021} supplies an integral local
pseudodeformation action even in scalar residual blocks. This resolves
a local representation-theoretic obstacle, but leaves a global one:
a trace does not record a residual extension class, and globally
reducible families can be too large for a direct connectedness argument.
We work with actual representations and temporarily record local
invariant lines to obtain the necessary dimension loss. These auxiliary
marks are used in connectedness; the propagation theorem imposes no
ordinary condition at the dyadic places.

\subsection{From occurrence to classicality and descent}

A \emph{pro-modular locus} is a closed trace locus in a Galois
pseudodeformation spectrum defined by a completed Hecke algebra of
definite quaternionic forms. No local condition is imposed at~\(2\).
Its characteristic-zero points give Hecke eigensystems in completed
forms. The central problem is to pass from this occurrence statement
to a classical cuspidal form with the prescribed Galois representation.
We do so in three stages, summarized in Figure~\ref{fig:proof-route}.

\paragraph{Potential occurrence for a whole family.}
The first stage places the target and the deformation families needed
later in one pro-modular trace locus over a common solvable totally
real extension \(F\), completely split at~\(2\). Two constructions
work together. Localized propagation starts at a characteristic-two
curve whose generic representation remains absolutely irreducible
under finite restriction. If the curve is pro-modular and has finite
local images at the allowed non-dyadic primes, patching gives
potential pro-modularity of every irreducible deformation locus
containing it (Theorem~\ref{thm:localized-propagation}). The local
deformation problem at~\(2\) is unrestricted.

To produce the curves to which this result applies, we construct
abelian \(2\)-extensions \(\mathbb Q\subset A\subset F\), keeping the
number of allowed ramified primes under control while their degrees
grow. In the scalar and reducible cases we mark a small proportion of
the dyadic places. Recording an invariant line at those places retains
information that traces alone would lose. We select the open part of
the residual fiber on which the marked lines are transverse to every
globally invariant line. A Selmer calculation makes the globally
reducible locus too small to disconnect these charts. A modular seed
and a curve cut from a component containing the target enter the same
connected residual fiber. Applying propagation along it yields
Theorem~\ref{thm:potential-promodularity}. The conclusion concerns the
whole closed trace image of the prescribed family at one common tame
level; it does not require that this family itself have the dimension
of the larger component used to find the curve.

\paragraph{Classical vectors at the target.}
When the local representation at~\(2\) is irreducible, the local
Langlands correspondence supplies locally algebraic vectors. For a
reducible local target, we choose an invariant line with the larger
cyclotomic exponent and retain its character as a parameter. The
resulting family records trace and character jointly and is finite
over a one-dimensional weight space. Its continuous members are not
all de Rham; the useful points are a dense set of specially chosen
weights. After a global twist the corresponding representations have
cyclotomic exponents \(\{1,0\}\), or Hodge--Tate weights
\(\{-1,0\}\). Their \emph{duals} have Hodge--Tate weights
\(\{0,1\}\) and satisfy the ordinary potentially crystalline
hypotheses of Thorne's theorem.

Those classical points determine the order of the two characters in
the ordinary induction. A finite closed joint support transports this
order to the target. A quotient, lattice and Tate-module construction
then supplies a genuine eigenvector in the target eigenspace. Finally,
a direct polynomial-orbit argument turns its character and exponent
inequality into a nonzero locally algebraic vector. This last step
avoids assuming that algebraic vectors commute with specialization.

\paragraph{Descent of a classical packet.}
The vector gives a regular cuspidal automorphic representation over
\(F\). We descend along prime-degree cyclic steps in a solvable tower.
At each step, a finite-character discrepancy between the descended
Galois representation and the target is removed by a Hecke-character
twist. Descent is applied only after classicality has been established;
no descent statement for pro-modular loci is needed.

\begin{figure}[tp]
\centering
\begin{tikzpicture}[
 box/.style={draw,rounded corners=2pt,align=center,font=\small,
 text width=5.7cm,inner sep=7pt},
 arr/.style={-{Latex[length=2mm]},thick},node distance=9mm]
\node[box] (prop) {Unrestricted propagation at a curve\\
Sections~\ref{sec:deformation}--\ref{sec:propagation}};
\node[box,right=7mm of prop] (conn) {Controlled fields and transverse connectedness\\
Sections~\ref{sec:fields}--\ref{sec:connectedness}};
\draw[arr] (prop.east) -- (conn.west);
\node[box,below=14mm of prop,xshift=3.2cm,text width=8cm] (occ)
{Potential pro-modularity of the target family\\over one common field \(F\)};
\draw[arr] (prop.south) -- (occ.north west);
\draw[arr] (conn.south) -- (occ.north east);
\node[box,below=13mm of occ,xshift=-3.2cm] (irr)
{Locally irreducible target:\\local Langlands correspondence};
\node[box,below=13mm of occ,xshift=3.2cm] (red)
{Locally reducible target:\\finite ordinary support,\\dense classical points\\and polynomial orbit};
\draw[arr] (occ.south) -- (irr.north);
\draw[arr] (occ.south) -- (red.north);
\node[box,below=14mm of irr,xshift=3.2cm,text width=8cm] (vec)
{Locally algebraic target vector\\and cuspidal packet over \(F\)\\
Section~\ref{sec:classicality}};
\draw[arr] (irr.south) -- (vec.north west);
\draw[arr] (red.south) -- (vec.north east);
\node[box,below=of vec,text width=8cm] (end)
{Prime-degree cyclic descent and finite-character corrections\\
Classical cuspidal eigenform over \(\mathbb Q\)\\up to Tate twist};
\draw[arr] (vec) -- (end);
\end{tikzpicture}
\caption{The proof route. The connectedness construction uses the
propagation theorem; the two upper boxes are cooperating inputs to
occurrence. The target enters through that construction. The common
field \(F\) is solvable, totally real, of even degree and completely
split at~\(2\). The ordinary family is used in the reducible local
branch. Both branches produce a classical packet before descent.}
\label{fig:proof-route}
\end{figure}

\subsection{The constructions used in the proof}

Several established methods must be kept compatible at the
characteristic-two curve. B\"ockle--Iyengar--Pa\v{s}k\=unas develop
local deformation theory with Laurent-series residue fields and
coefficient orders, including the auxiliary-variable comparison of
completions \cite[Section~3.5]{BIP}. We use this framework
with an extra coefficient variable and global representation and
incidence spaces. Lemma~\ref{lem:curve-completion} proves the resulting
identification of continuous deformation and obstruction problems,
including the topology on the residue field and the dimension change.

The propagation construction follows the auxiliary-level method of
Taylor--Wiles, Kisin's formulation relative to local deformation rings,
and the characteristic-two determinant-twist argument of
Khare--Wintenberger
\cite{TaylorWiles1995,Kisin2009FiniteFlat,KhareWintenberger2009II}.
Pan's localization at one-dimensional primes is the direct global
precedent \cite{Pan2022}; Zhang's potential formulation is used in
the same overall progression \cite{Zhang2025}. At the Laurent-series
residue field we patch finite jets, meaning quotients by fixed powers
of the local maximal ideal. The exactness argument works with
filtrations whose successive quotients are fixed old-level modules.
A uniform Artin--Rees bound preserves their actual old deformation
actions. Determinant twists and explicit quadratic sign intertwiners
then give support on every local component. All of these action
comparisons are retained integrally.

The marked-line construction serves a different purpose. Its local
factors are integral even over the characteristic-two real-place
base, and transversality gives a strict Selmer loss at every mark.
The number of marks grows more slowly than the total degree, so this
loss is compatible with the global dimension estimates. The argument
uses the unrestricted propagation theorem after entering a transverse
chart; it does not replace that theorem by ordinary patching.

\paragraph{Organization.}
Section~\ref{sec:preliminaries} fixes conventions and removes the
virtually solvable case. Section~\ref{sec:deformation} constructs the
representation spaces, their curve completions and the local models.
Section~\ref{sec:completed} compares the actual Galois and local
block actions on completed forms and proves the Hecke-dimension bound.
Section~\ref{sec:propagation} proves localized propagation.
Section~\ref{sec:fields} supplies the field and Selmer estimates, and
Section~\ref{sec:connectedness} combines them with a modular seed to
obtain potential pro-modularity. Section~\ref{sec:classicality}
constructs classical vectors and carries out solvable descent.
